\section{Conclusion}

The assignment's advantage for Eq.~4 over Eq.~3 does not hold up. About half of all draws of the same pipeline (\N{rep3:ge3} of 100) put Eq.~4 ahead on three or more data sets, but the interval of the mean difference excludes zero only on Seeds (for Eq.~4) and Iris (for Eq.~3), unadjusted, on the report's one split, and with kernels chosen on the test set; none of the tuned study's comparisons of Eq.~4 with Eq.~3 survives Holm's correction. In paired differences pooled over draws, target sizes, and the data sets other than Iris on which each family was run, Eq.~3, Eq.~4, and the very sparse matrices stay within about a point of a dense Gaussian projection (balanced accuracy for Glass; across all five data sets and the four non-Gaussian families, only \N{z:equiv_vs_gauss} of the \N{z:n_vs_gauss} primary comparisons with it are equivalent within $\pm 2$ points, mostly at larger $k$); where $d$ is tiny and sparsity drops whole features and can leave a column of the matrix all zero, as on Iris, Eq.~4 trails by more (the loss comes from the draws with a zero column), and one nonzero per column loses everywhere. When the data can be fit, PCA is a better choice at small $k$. If the map must be data-independent, Achlioptas's sparse matrix is a reasonable substitute for a dense one here, and $k$ must be large for the mean to come within 2 points of all features for the Gaussian, Eq.~3, and Eq.~4: 12 of \N{z:d:breast_cancer} features on Breast Cancer, 4 to 6 of \N{z:d:seeds} on Seeds, depending on the family, and on Wine only 12 of \N{z:d:wine}, where the intervals still reach \N{mag:pd:wine:12:gauss:full:lo} points below. No target size here meets the Johnson--Lindenstrauss sufficient condition.
